% ====================================================================
% ch1_sec16b_lcoomega.tex — [LCO-Ω] LCO per-peak Ω(정칙용액 커널) + 전자항 on/off 토글
%   삽입: ch2 마스터 — §16 LCO 세 peak(ch1_sec16_lcopeak) 뒤, §17 MSMR 앞.
%   v1.0.24 신규(@3 per-peak Ω·#7 문구 정정·전자항 토글).
%   저자 = comp_R1 W9 초안(master 통합판).
%   ★토글 기본값 = include_electronic_entropy=False (OFF; 사용자 사양 "커브 구할 땐 빼고"·브리프 "기본 OFF").
%     상온 커브는 토글과 무관하게 전자항 무영향(U(T_ref) 불변, ΔH^eff 재기준).
%     True(옵션) = v1.0.19/v1.0.23 상시-ON 거동 재현 --- G1 회귀는 이 ON 경로로 bit-exact 검증.
% ====================================================================
\subsection{LCO per-peak $\Omega$ 커널과 전자항 on/off 옵션 --- 평균장 축약의 지위}\label{ssec:lco-omega}
LCO 하프셀 dQ/dV 는 흑연과 \emph{같은} MSMR(전하 보존 반전 $\to$ 평형 peak, 식~\eqref{eq:lco-eqpeak})로 닫히고,
전이 파라미터만 양극 값으로 바뀐다(표~\ref{tab:lco-staging}). 이 소절은 그 위에 두 가지를 정합하게 얹는다:
(i) feature 별 상성격 차이를 담는 \emph{per-peak $\Omega_j^\cat$ 커널}과 그와 얽힌 \S\ref{sec:lco-hys} 문구의
\#7 정정, (ii) 전자 엔트로피 항을 \emph{커브에는 무영향$\cdot$가역열에만 작용}하는 on/off 토글로 재배치하는 규칙.
새 물리는 없다 --- 식~\eqref{eq:lco-gxi}$\cdot$\eqref{eq:lco-gpp}$\cdot$\eqref{eq:lco-spinodal} 의 정칙용액 골격과
식~\eqref{eq:dSegate}$\cdot$\eqref{eq:U1T2} 의 전자항을 \emph{어느 슬롯이 무엇을 담는가}로 정리할 뿐이다.

\textbf{(a) 출발 --- feature 별 상성격은 per-peak $\Omega_j^\cat$ 가 가른다.} LCO 세 전이(표~\ref{tab:lco-staging})는
상성격이 균일하지 않다: $\sim$3.90 V 의 MIT(절연체$\to$금속) 계열은 회절이 두 상 공존을 보이는 \emph{두-상 sharp}
전이이고\cite{reimers1992,menetrier1999}, $x\approx\tfrac12$ 부근 order--disorder 계열은 정렬 상호작용이 주도하는
전이다\cite{vanderven1998}. 이 차이를 forward 모델은 전이별로 다른 $\Omega_j^\cat$ 하나로 담는다 --- 흑연
판정자와 같은 문턱에서(식~\eqref{eq:lco-gpp}$\cdot$\eqref{eq:lco-spinodal}), 자유에너지 곡률이 정하는 정칙용액
(Frumkin) 단일 peak 커널로 닫힌다:
\begin{equation}
\frac{\partial^2 g_j^\cat}{\partial\xi_j^2}=\frac{RT}{\xi_j(1-\xi_j)}-2\Omega_j^\cat
\ \Longrightarrow\
\boxed{\;\Big(\frac{\dd Q}{\dd V}\Big)^\eq_j=\frac{Q_j^\cat\,F}{\big|\,RT/[\xi_j(1-\xi_j)]-2\Omega_j^\cat\,\big|}\;}
\ \xrightarrow[\ ]{\ \Omega_j^\cat\to0\ }\
Q_j^\cat\,\frac{\xi_{\eq,j}(1-\xi_{\eq,j})}{w_j^{\eq}}\ \ (w_j^{\eq}\equiv RT/F),
\label{eq:lcoomega-kernel}
\end{equation}
곧 흑연 판정자(\S\ref{ssec:gr-2L} 식~\eqref{eq:gr2l-disc})$\cdot$Si-host Frumkin 종과 \emph{같은} 분모
$RT/[\xi(1-\xi)]-2\Omega$ 를 쓰는 한 커널이다 --- 이 분모는 자유에너지 곡률 $\partial^2 g_j^\cat/\partial\xi_j^2$
[J/mol](식~\eqref{eq:lco-gpp})이고, 그 역수에 $F$ 를 실은 것이 $\dd\xi/\dd V$ 다. \emph{표기 정합(진행좌표는 $\xi$
하나).} 박스의 $\xi_j$ 는 전위 $V$ 에서의 평형 진행률 $\xi_{\eq,j}(V)$(음함수 $V{=}V_{\eq,j}$ 의 근,
식~\eqref{eq:lco-Veq})로 읽어 극한식의 $\xi_{\eq,j}$ 와 \emph{같은 한 좌표}이며 --- 커브 전개를 $\xi$ 로 일관되게
쓰고 새 기호 $\theta$ 를 들이지 않으므로 $\xi$-$\theta$ 혼선이 없다.

\emph{$\Omega_j^\cat\to0$ 환원(명시$\cdot$자기정합).} 두 몫이 동시에 닫힌다. (i) \emph{폭.} 곡률이
$RT/[\xi(1-\xi)]$ 로 남아 $\dd\xi_{\eq,j}/\dd V=F\,\xi(1-\xi)/RT$ 이므로 커널은
$Q_j^\cat F\,\xi_{\eq,j}(1-\xi_{\eq,j})/RT=Q_j^\cat\,\xi_{\eq,j}(1-\xi_{\eq,j})/w_j^{\eq}$, 곧 \emph{이상용액
Nernst 폭} $w_j^{\eq}{\equiv}RT/F$ 로 닫히고, 이는 현상학 폭 $w_j{=}n_jRT/F$(식~\eqref{eq:lco-eqpeak} 의 폭
서식)의 \emph{$n_j{=}1$ 지점}이다. (ii) \emph{중심.} 동시에 $\Omega_j^\cat\!\to\!0$ 이면 히스 gap 이 사라져
($\Delta U_j^{\hys,\cat}{=}0$, 식~\eqref{eq:lco-dUhys}) 분기 중심이 $U_j^{\,d,\cat}{=}U_j^\cat$ 로 접히므로,
$\xi_{\eq,j}$ 는 중심 $U_j^\cat\cdot$폭 $RT/F$ 의 \emph{순 로지스틱}이 되어 평형 peak 식~\eqref{eq:lco-eqpeak} 의
\emph{이상 핵}과 정확히 겹친다. 따라서 $\Omega$ 미지정($=0$)은 식~\eqref{eq:lco-eqpeak} 경로로 bit-exact
폴백한다(커널 분기를 타지 않음). \emph{폭 슬롯 분리(이중계산 방지).} 환원의 $w_j^{\eq}{=}RT/F$ 는 커널이
예측하는 \emph{열역학 이상} 폭이고, 실제 LCO 세 전이가 놓이는 $\Omega_j^\cat>2RT$ 두-상 측의 \emph{관측} 폭은
이 극한이 아니라 현상학적 자유 폭 $w_j{=}n_jRT/F$($n_j\neq1$, \S\ref{sec:lco-peak} ``폭의 지위'')가 담는다 ---
커널의 두-상 델타가 관측 폭을 $w_j$ 로 넘기므로 $w_j^{\eq}$(이상 핵)와 $w_j$(두-상 broadening)는 서로 다른
층으로 갈려 분모 $|RT/[\xi(1-\xi)]-2\Omega|$ 와 $w_j$ 가 충돌 없이 각자의 역할을 맡는다.

각 전이의 sharp/broad 는 \emph{피팅된} $\Omega_j^\cat$ 이 $2RT$ 를 넘는지가 정한다(\S\ref{sec:lco-hys}
two-phase calibration) --- MIT 계열은 $\Omega_j^\cat>2RT$ 로 두-상 델타에 가깝고(관측 폭은 위 $w_j$ 가 담음),
도핑 시료의 order--disorder 계열은 $\Omega_j^\cat$ 이 $2RT$ 아래로 억제되면(식~\eqref{eq:lco-dope}) 상대적으로
\emph{broad} 하다. 곧 하나의 per-peak $\Omega_j^\cat$ 축이 세 feature 를 연속적으로 sharp$\leftrightarrow$broad
로 조율하며, 이 자유도가 LCO 곡선을 흑연 균일 폭보다 잘 맞춘다(정칙용액 커널 개선 $+1.25\%$p, 우리 ablation 진단).

\textbf{(b) 연산 --- \#7 정정: $\Omega_j^\cat$ 는 평균장 축약이지 미시 질서상 안정성 주장이 아니다.}
\S\ref{sec:lco-hys} 가 order--disorder 를 정칙용액으로 접을 때 ``$\Omega_j^\cat>2RT\Leftrightarrow x\!\approx\!\tfrac12$
질서상 안정''(식~\eqref{eq:br-vanderven1998-1})이라는 등가 문구는 \emph{과대 해석의 소지}가 있어 다음으로
정정한다 --- 식$\cdot$값$\cdot$슬롯은 그대로 두고 \emph{읽기(문구)}만 정확화한다(물리 불변):
\begin{equation}
\boxed{\;
\Omega_j^\cat\ \equiv\ \text{전이 $j$ 진행좌표 $\xi_j$ 의 \emph{유효 평균장 쌍상호작용 축약}}\ \ (\text{미시 질서구조가 아님}),
\quad
\Delta S_j^\config\ \perp\ \Omega_j^\cat\;}
\label{eq:lcoomega-hash7}
\end{equation}
--- 즉 $\Omega_j^\cat$ 은 Van der Ven 등\cite{vanderven1998} 의 다체 클러스터 전개를 전이 창의 진행률 한 좌표로
접었을 때 남는 \emph{최근접 쌍 몫의 평균장 값}이고, 그것이 $2RT$ 를 넘는다는 것은 그 유효 좌표의 볼록성 상실
(두-상 문턱)을 뜻할 뿐, $x{=}\tfrac12$ Li$\cdot$빈자리 초격자의 \emph{미시 질서구조}를 $\Omega_j^\cat$ 이 기술한다는
말이 아니다. 곧 식~\eqref{eq:br-vanderven1998-1} 자체(값$\cdot$구조)는 손대지 않고, 그 ``$\Leftrightarrow$ 질서상
안정'' 읽기를 ``유효 좌표 볼록성 상실 $=$ 두-상 문턱''으로 다시 읽을 뿐이다. 질서화가 나르는 \emph{배치
엔트로피}(config)는 별도 슬롯 $\Delta S_j^\config$(식~\eqref{eq:lco-decomp})에 들어가며 상호작용 에너지
$\Omega_j^\cat$[J/mol]와 차원부터 \emph{직교}한다[J/(mol\,K)] --- 그래서 정렬의 charge-order $\Delta S$
(값$\cdot$배정 양쪽 tier C)를 ``$\to\Omega_j^\cat$''로 직접 연결하지 않는다(\S\ref{sec:lco-hys-od} 혼동 가드
계승 --- 이 정정은 그 가드를 이 소절에 이어 붙인 것이지 새 물리가 아니다). 이 정정이 \emph{gap 슬롯}
($\Omega_j^\cat$)과 \emph{엔트로피 슬롯}($\Delta S_j^\config$)의 이중계산을 문구 수준에서 막는다.

\textbf{(c) 중간식 --- 전자항은 $T_\mathrm{ref}$ 동결로 상온 커브에 무영향, $\partial U/\partial T$ 에만 작용.}
LCO 고유의 전자 엔트로피 $\Delta S_{e,j}(x,T)\propto T$(MIT 게이트, 식~\eqref{eq:dSegate})는 반응 엔트로피 슬롯에
들어가되, 그 작용처가 커브와 가역열에서 다르다. 기준온도 $T_\mathrm{ref}$ 에서 중심 전위는 식~\eqref{eq:Uj} 로
\begin{equation}
U_j(T_\mathrm{ref})=\frac{-\Delta H_{\rxn,j}^\mathrm{eff}+T_\mathrm{ref}\,\Delta S_{\rxn,j}^{\cat}(T_\mathrm{ref})}{F},
\qquad
\Delta H_{\rxn,j}^\mathrm{eff}\equiv\Delta H_{\rxn,j}-T_\mathrm{ref}\,\Delta S_{e,j}^{\,\mathrm{mol}}(T_\mathrm{ref})\ \ (\text{재기준}),
\label{eq:lcoomega-Tref}
\end{equation}
곧 $T_\mathrm{ref}$ 에서 전자항의 상수 오프셋 $T_\mathrm{ref}\Delta S_{e,j}$ 는 유효 엔탈피 $\Delta H^\mathrm{eff}$
로 흡수되어 $U_j(T_\mathrm{ref})$ 를 \emph{불변}으로 둔다. 따라서 상온($T{=}T_\mathrm{ref}$) 단일 커브는 전자항을
켜든 끄든 같다 --- 전자항이 커브에 남기는 것은 $U_j(T_\mathrm{ref})$ 가 아니라 오직 그 \emph{온도 기울기}
$\partial U_j/\partial T=\Delta S_{\rxn,j}^\cat/F$ 뿐이고, 이는 다온도 데이터(가역열, Chapter 1 Part T
\S\ref{sec:revheat})에서만 관측된다. 실측 뒷받침 --- LCO plain MSMR 적합이 $R^2{=}0.944$ 로 흑연 $0.940$ 과
사실상 같아(우리 진단), 상온 커브 재현에 전자항이 불필요함을 보인다.

\textbf{(d) 박스 --- 전자항 on/off 옵션(기본 off).} 이를 옵션 하나로 고정한다:
\begin{equation}
\boxed{\;\text{전자항 옵션}=
\begin{cases}
\text{off (기본)} & \Delta S_{e}\ \text{를 }\partial U/\partial T\ \text{에서 제외}\ \to\ U_j(T_\mathrm{ref})\ \text{보존(상온 커브 불변)},\\[2pt]
\text{on (옵션)} & \partial U_j/\partial T\ \text{에 }\Delta S_{e,j}(T_\mathrm{ref})\ \text{상수 오프셋 활성}\ \to\ \text{가역열 정량}.
\end{cases}\;}
\label{eq:lcoomega-toggle}
\end{equation}
규칙의 요지 --- \emph{상온 커브(dQ/dV @$T_\mathrm{ref}$)는 토글과 무관하게 전자항 무영향}이다($U(T_\mathrm{ref})$
불변, 식~\eqref{eq:lcoomega-Tref}). 기본을 \textbf{off}(전자항 제외)로 두어도 상온 커브는 전자항과 무관하므로
곡선 산출에서 전자항 값의 선택은 무의미하다. 토글이 실제로 가르는 것은 오직 $\partial U/\partial T$(가역열$\cdot$다온도)다
--- \textbf{on 옵션}은 전자항의 \emph{가역열 기여}를 켜서 정량화하고 최종 유지/제거를 판정하는 ``있고 없고 선택''
장치다. \textbf{on} 경로는 전자항을 항상 포함하는 거동을 bit-exact 로 보존하고, 기본 \textbf{off} 는
$\Delta H^\mathrm{eff}$ 재기준으로 $U(T_\mathrm{ref})$ 를 불변으로 유지한다.

\medskip
\textbf{세 쟁점의 정리.} \emph{(i) 전자항의 작용 범위.} 전자항은 상온 커브에는 불필요하고 가역열에는 필요하다 ---
상온 dQ/dV 는 식~\eqref{eq:lcoomega-Tref} 로 전자항과 무관(plain MSMR $R^2{=}0.944\!\approx\!$흑연
$0.940$)하나, MIT 구간의 $\partial U/\partial T$ 온도 서명은 전자항 없이 config 단독으로 재현되지 않는다
(\S\ref{sec:lco-electronic}: config 전용 계산\cite{ml2024} 도 order--disorder 소관, 전자 자유도 미포함) --- ``작지만
빠지면 MIT 를 못 그린다''는 서명이 \emph{가역열 축}에 나타난다. 곧 토글은 두 축을 분리하는 장치다.
\emph{(ii) per-peak $\Omega$ 의 지위.} $\Omega_j^\cat$ 은 자유 봉우리 추가가 아니라 회절
\cite{reimers1992}$\cdot$제일원리 클러스터 전개\cite{vanderven1998}$\cdot$전자 상도표\cite{motohashi2009} 가
받치는 \emph{상성격} 축이며, 판정자는 물리적 문턱($\Omega_j^\cat\!\lessgtr\!2RT$)이지 임의 폭이 아니다. 개선폭
$+1.25\%$p 는 \emph{완만}하고 물리 동기가 분명하다 --- 봉우리 수를 늘리는 대신 $\Omega$ 한 축을 조율하므로 과적합이 아니다.
\emph{(iii) \#7 정정의 범위.} \#7 정정은 식$\cdot$값$\cdot$슬롯 배치를 전부 유지하며 $\Omega_j^\cat$ 를
\emph{평균장 축약}으로 규정하므로(식~\eqref{eq:lcoomega-hash7}; 식~\eqref{eq:br-vanderven1998-1} 원식 불변) 물리
내용은 바뀌지 않는다.

\begin{warnbox}
\textbf{정직한 한계.} (i) \emph{O3-LCO 온도 의존은 다온도 데이터 없이 미검증(tier 없음 --- round-trip 미확정).}
전자항의 유일한 관측 신호는 $\partial U_1/\partial T$ 가 $T$ 에 선형으로 내려가는 것(U-이동 $\propto T^2$,
식~\eqref{eq:U1T2})인데, 상온 단일 커브로는 이를 볼 수 없다 --- 회사 다온도 셀 데이터가 들어와야
$g_{\max}\cdot x_\mathrm{MIT}\cdot\Delta x_\mathrm{MIT}$ 게이트 세 초기값과 $\Delta S_e$ 부호$\cdot$기울기가
round-trip 으로 확정된다(현 상태: 연속 $g(E_F,x)$ 는 1차 문헌 부재 갭 G2, \S\ref{sec:lco-Se}). (ii) \emph{전자항
크기의 tier 분리.} 함수형 $S_e{=}\tfrac{\pi^2}{3}k_B^2Tg(E_F)$ 는 tier A(Sommerfeld 표준), 완전 metal 끝점
$g(E_F)\!\approx\!13$ e/eV/atom 도 그 한 점에서 tier A\cite{motohashi2009}, 그러나 MIT 를 가로지르는 \emph{연속
곡선}은 tier 없음(모델 게이트 가정) --- 골 깊이 $\sim-45.7$ J/(mol\,K)($T_\mathrm{ref}$ 동결 시연값)는 시연
감각이지 실측 예측이 아니다. (iii) \emph{6+ feature 는 폐기.} 흑연과 같은 규율 --- LCO 도 하프셀 상한
($\le$4.2--4.5 V) 안에서 세 feature 가 실측 상한(고전압 O3$\to$H1-3 는 범위 밖 옵션)이며, 이를 넘어 봉우리를
쪼개는 것은 상도표 미지원 curve-fitting 이다. (iv) \emph{$\Delta S_j^\config$ 수치는 값$\cdot$배정 양쪽 tier C}
(원전 재확인 중, \S\ref{sec:lco-hys-od} 혼동 가드) --- \#7 정정은 이 수치를 $\Omega$ 로 옮기지 말라는 경계이지
수치 확정이 아니다.
\end{warnbox}

\begin{keybox}
\textbf{LCO $\Omega$$\cdot$전자항 다리.} feature 별 상성격은 per-peak $\Omega_j^\cat$(식~\eqref{eq:lcoomega-kernel})가
피팅된 $2RT$ 초과 여부로 sharp/broad 를 가르며($\Omega\!\to\!0$ 에서 $w_j^{\eq}{=}RT/F$ 로 식~\eqref{eq:lco-eqpeak}
이상 핵에 bit-exact 폴백 --- 두-상 관측 폭은 현상학 $w_j{=}n_jRT/F$ 가 별도 층으로 담음), \#7 정정으로
$\Omega_j^\cat$ 은 \emph{평균장 축약}(미시 질서상 아님)$\cdot$config 엔트로피는 별도 슬롯 직교
(식~\eqref{eq:lcoomega-hash7})다. 전자항은 $T_\mathrm{ref}$ 동결로 상온 커브에 무영향($\Delta H^\mathrm{eff}$
흡수, plain MSMR $R^2{=}0.944\!\approx\!$흑연 $0.940$)이고 $\partial U/\partial T$(가역열)에만 작용하므로
전자항 옵션은 \emph{기본 off}(전자항 제외$\cdot$상온 커브 불변)$\cdot$\emph{on 옵션}
(가역열에 전자항 활성)으로 둔다(식~\eqref{eq:lcoomega-toggle}).
\emph{정직 한계}: O3-LCO 온도 의존은 다온도 데이터 없이 미검증, 연속 $g(E_F,x)$ 는 tier 없는 모델 게이트.
\end{keybox}

% 자산(comp_R1 W9 → master 통합·refine_b 정정): [LCOΩ-1 per-peak Ω_j^cat 커널 eq:lcoomega-kernel(+1.25%p);
%   ★Ω→0 환원 명시: w_j^eq=RT/F(=w_j=n_jRT/F 의 n_j=1)·gap=0(eq:lco-dUhys)→중심 U_j·순 로지스틱=eq:lco-eqpeak 이상 핵 bit-exact 폴백; ξ 단일 좌표(θ 없음)·두-상 폭은 현상학 w_j 별도 층]
%   [LCOΩ-2 #7 정정 boxed eq:lcoomega-hash7(Ω=평균장 축약·config⊥Ω; eq:br-vanderven1998-1 원식 불변·읽기만 정정·sec:lco-hys-od 가드 계승)]
%   [LCOΩ-3 전자항 T_ref 동결 재기준 eq:lcoomega-Tref(커브 무영향·R²=0.944≈0.940)]
%   [LCOΩ-4 include_electronic_entropy 토글 boxed eq:lcoomega-toggle(★기본 False=OFF=전자항 제외·상온커브 불변·True 옵션=가역열 활성=v1.0.19 거동·G1 은 ON 경로 검증)]
%   [LCOΩ-5 적대 3문+정직 한계 warnbox(전자항 커브불필요/가역열필요·O3 다온도 미검증 tier없음·6+ 폐기·config tier C)]
